\subsection*{Synchronization in Practice: Pipelines, Buffers, and Events}

Section~\ref{sec:3-8} introduced three control-flow primitives: \textbf{\texttt{if}} for predicated execution, \textbf{\texttt{inthreads.\allowbreak{}async}} for structured concurrent regions, and \textbf{events} (\texttt{shared event} / \texttt{wait} / \texttt{trigger}) for inter-region signaling. This part of the guide puts them to work. You will see two progressively complex kernel patterns, each demonstrating how the primitives compose to solve real synchronization problems.

The first pattern --- \textbf{double buffering with \texttt{swap}} --- uses a single thread group that interleaves loading and computing within one program. No \texttt{inthreads.\allowbreak{}async}, no events: just two buffer handles and a rotation. The second pattern --- \textbf{the full 1P1C event pipeline} --- splits loading and computing into separate concurrent programs with \texttt{inthreads.\allowbreak{}async} and coordinates them with event arrays.

Both patterns solve the same underlying problem: \textbf{you cannot read a buffer while someone is writing to it}. They differ in how they structure the solution.
